\section{Refining the phase-ball cover}\label{sec:refinement}

The purpose of Stage B is to replace a finite collection of certified
phase balls by balls of a smaller prescribed radius.  The centers need
not be exact Hadamard matrices, and the argument makes no regularity
assumption on the Hadamard locus.  In particular, the tangent coordinates
used below do not parametrize that locus: they organize an exhaustive
search, while the normal estimates from Section~\ref{sec:phase-cover}
control the remaining directions.

All vector spaces of phases in this section are real.  A phase matrix is
identified with a vector in $\R^{36}$, and $^*$ denotes transpose on real
spaces.  Write $\mathcal V=\operatorname{im}\Pi$ for the horizontal
subspace.  For a certified center $P$ and a radius $R\geq0$, let
\begin{equation}\label{refine:solution-set}
 \mathcal S(P,R)=
 \bigl\{[e(P+\delta)]:\delta\in\mathcal V,\quad
       \|\delta\|\leq R,\ F(P+\delta)=0\bigr\}.
\end{equation}
Here brackets denote the equivalence classes used in the preceding
sections, including the permitted signed permutation and transpose
operations.  These equivalences preserve the extendibility property
being tested.  Thus a cover by the sets \eqref{refine:solution-set} is
sufficient even when different parts of it use different arrangements
of the same matrix.

\subsection{A certified relation between two centers}

For the Root at $P$, denote its stored tangent matrix by
$T:\R^k\longrightarrow\R^{36}$ and put
\begin{equation}\label{refine:constants}
 \begin{split}
 o_P&=\mathtt{ort}_P,\qquad
 \alpha_P=\mathtt{id}_P,\qquad \gamma_P=\mathtt C_P,\\
 \mu_P&=\mathtt{mnorm}_P,\qquad
 \varepsilon_P=\mathtt{froot}_P,\qquad
 n_P(s)=\mathtt{n0}_P+\alpha_Ps+\gamma_Ps^2.
 \end{split}
\end{equation}
All these constants are nonnegative.  We use the same notation, with
subscript $Q$, for a second certified Root with center $Q$ and tangent
matrix $S:\R^{k_Q}\longrightarrow\R^{36}$.  In particular,
\begin{equation}\label{refine:gram-bounds}
 \|T^*T-I\|\leq o_P,
 \qquad \|S^*S-I\|\leq o_Q,
 \qquad \|F(Q)\|\leq\varepsilon_Q.
\end{equation}
We will apply the normal and difference inequalities
\eqref{phase:root-normal} and \eqref{phase:root-difference}.

The following elementary identity explains why approximate, rather
than exact, orthogonality is sufficient throughout the construction.

\begin{lemma}\label{refine:approximate-projection}
Put $A=I-TT^*$.  For every $x\in\R^{36}$, with $\zeta=T^*x$,
\begin{equation}\label{refine:projection-estimates}
 \|T\|\leq\sqrt{1+o_P},\qquad
 \|A\|\leq1+o_P,\qquad
 \|x\|^2\leq(1+o_P)\|\zeta\|^2+\|Ax\|^2.
\end{equation}
These statements do not require $x$ or the columns of $T$ to be
horizontal.
\end{lemma}

\begin{proof}
The Gram bound gives the first inequality.  The eigenvalues of $TT^*$
lie in $[0,1+o_P]$, so those of $I-TT^*$ lie in $[-o_P,1]$; hence its
norm is at most $\max(1,o_P)\leq1+o_P$.  Finally,
\[
 \|Ax\|^2
 =\|x\|^2-2\|\zeta\|^2+\|T\zeta\|^2
 \geq\|x\|^2-(1+o_P)\|\zeta\|^2,
\]
using $\|T\zeta\|^2\geq(1-o_P)\|\zeta\|^2$.
\end{proof}

Let $D_0$ be the unwrapped dephasing operator,
\[
 (D_0X)_{ij}=X_{ij}-X_{i0}-X_{0j}+X_{00}.
\]
Choose any integer matrix $N$ and define
\begin{equation}\label{refine:alignment}
 d_q=\Pi\bigl(D_0Q-D_0P-N\bigr),\qquad
 \eta_q=T^*d_q,\qquad n_q=\|A d_q\|.
\end{equation}
The implementation chooses $N$ by wrapping differences of cached
dephased entries.  The chosen lift need not minimize any distance.
Indeed, $D_0X-X\in\ker\Pi$, so
\[
 Q-P-d_q\in N+\ker\Pi.
\]
Vectors in $\ker\Pi$ are row and column rephasings.  Consequently
$P+d_q$ and $Q$ differ only by integer phases and rephasings.

Fix a parent radius $R$ and an exact solution displacement
$\delta\in\mathcal V$ with $\|\delta\|\leq R$.  Set
\begin{equation}\label{refine:relative-displacement}
 \eta=T^*\delta,\qquad
 h'=\delta-d_q,\qquad d=\|h'\|,\qquad
 \xi=T^*h'=\eta-\eta_q,\qquad \nu=Ah'.
\end{equation}
Both $d_q$ and $h'$ are horizontal.  Moreover $Q+h'$ is equivalent by
the same rephasing to $P+\delta$, and therefore is exact.  Thus the
Root estimates at both centers apply to this one solution.

\begin{proposition}\label{refine:relation-estimates}
Suppose $t\geq\|\xi\|$.  Choose upper bounds
\begin{equation}\label{refine:relation-constants}
 \begin{split}
 D&\geq R+\|d_q\|,\qquad
 v_0\geq n_P(R)+n_q,\\
 \lambda&\geq\alpha_P+\gamma_PD,\qquad
 e_0\geq\mu_P\varepsilon_Q.
 \end{split}
\end{equation}
Let $S_h=\Pi S$ and choose
\begin{equation}\label{refine:sigma-definition}
 \sigma\geq\|S_h^*A\|,
 \qquad
 \beta=(1+o_P)\sqrt{(1+o_P)(1+o_Q)}.
\end{equation}
Then
\begin{align}
 d&\leq D,
 & d^2&\leq(1+o_P)t^2+\|\nu\|^2,
 \label{refine:parent-distance}\\
 \|\nu\|&\leq
 \min\{v_0,(1+o_P)d,\lambda d+e_0\},
 \label{refine:normal-three}\\
 \|S^*h'\|&\leq\beta t+\sigma\|\nu\|,
 \label{refine:child-tangent}\\
 d^2&\leq(1+o_Q)\|S^*h'\|^2+n_Q(d)^2.
 \label{refine:child-distance}
\end{align}
\end{proposition}

\begin{proof}
The triangle inequality gives $d\leq R+\|d_q\|\leq D$.
Lemma~\ref{refine:approximate-projection}, applied to $h'$, gives the
other inequality in \eqref{refine:parent-distance}.  Also
\[
 \|\nu\|\leq\|A\delta\|+\|Ad_q\|
       \leq n_P(R)+n_q\leq v_0,
\]
because the coefficients of $n_P$ are nonnegative.  The operator bound
on $A$ gives the second candidate in \eqref{refine:normal-three}.
For the third, use \eqref{phase:root-difference} with its horizontal
argument $v=d_q$:
\[
 \|Ah'\|
 \leq\bigl[\alpha_P+\gamma_P(\|\delta\|+\|d_q\|)\bigr]d
       +\mu_P\|F(P+d_q)\|
 \leq\lambda d+e_0.
\]
Here residual-norm invariance under rephasing gives
$\|F(P+d_q)\|=\|F(Q)\|\leq\varepsilon_Q$.

To prove \eqref{refine:child-tangent}, first use horizontality of $h'$
to obtain $S^*h'=S_h^*h'$.  The identities
\[
 h'=T\xi+\nu,\qquad
 T^*\nu=(I-T^*T)\xi,\qquad
 \nu=A\nu+TT^*\nu
\]
therefore imply
\[
 S^*h'
 =S_h^*T\xi+S_h^*A\nu
       +S_h^*T(I-T^*T)\xi.
\]
Since $\Pi$ is an orthogonal projection,
\[
 \|S_h^*T\|
 \leq\|S_h\|\,\|T\|
 \leq\sqrt{(1+o_Q)(1+o_P)}.
\]
The first and third terms are thus bounded together by $\beta t$,
and the middle term by $\sigma\|\nu\|$.  This calculation also shows
why no assumption of exactly horizontal stored tangent columns, and
no omitted error term involving $\mathtt{hor}$, is hidden in the
estimate.  Finally apply the squared inequality at the child Root
and its normal bound to the horizontal exact-solution displacement
$h'$.  This gives \eqref{refine:child-distance}.
\end{proof}

The coefficient $\sigma$ has a direct small-matrix certificate.  With
$W=(AS_h)^*$, its row Gram is
\begin{equation}\label{refine:sigma-gram}
 G=WW^*=S_h^*A^2S_h,
 \qquad \|S_h^*A\|^2=\|G\|.
\end{equation}
The code constructs these rows by projecting each child tangent column
and subtracting its components against the parent columns, in that
order.  A certified upper bound for the norm of $G$, followed by a
square root, gives \eqref{refine:sigma-definition}.  For any exact real
symmetric $G$, either of the quantities
\[
 \max_i\sum_j|G_{ij}|,
 \qquad \sqrt{\|G^2\|_F}
\]
bounds $\|G\|$: the first is the symmetric absolute-row bound, while
$\|G\|^2=\|G^2\|\leq\|G^2\|_F$ proves the second.  The routine
\texttt{spectralUpper} uses their minimum with margins for arithmetic
and Gram perturbation, detailed in Section~\ref{sec:rounding}.
When $k_Q=0$, the operator is zero and the code sets $\sigma=0$.

\subsection{Distance iteration and the searched tangent radius}

The relation estimates yield a useful upper iteration without requiring
either contraction or a fixed-point theorem.  For fixed $t$, write
\begin{equation}\label{refine:iteration-functions}
 \begin{split}
 z&=\sqrt{1+o_P}\,t,\\
 v(b)&=\min\{v_0,(1+o_P)b,\lambda b+e_0\},\\
 \psi_P(b)&=\sqrt{z^2+v(b)^2},\\
 \psi_Q(b)&=
 \sqrt{(1+o_Q)(\beta t+\sigma v(b))^2+n_Q(b)^2}.
 \end{split}
\end{equation}
Start at
\begin{equation}\label{refine:iteration-start}
 b_0=\min\{D,\sqrt{z^2+v_0^2}\}
\end{equation}
and replace $b_j$ by
\begin{equation}\label{refine:iteration-step}
 b_{j+1}=\min\{b_j,\psi_P(b_j),\psi_Q(b_j)\}.
\end{equation}

\begin{lemma}\label{refine:iteration-certificate}
Every iterate in \eqref{refine:iteration-start}--\eqref{refine:iteration-step}
is an upper bound for the distance $d$ of every parent solution with
$\|\eta-\eta_q\|\leq t$.  The same conclusion holds if any proposed
bound is replaced by a larger certified number, or if the process is
stopped after any finite number of steps.
\end{lemma}

\begin{proof}
Both candidates for $b_0$ are upper bounds by
Proposition~\ref{refine:relation-estimates}.  If $d\leq b_j$,
\eqref{refine:normal-three} and nonnegativity of all coefficients give
$\|\nu\|\leq v(b_j)$.  Thus \eqref{refine:parent-distance} gives
$d\leq\psi_P(b_j)$.  Equations
\eqref{refine:child-tangent}--\eqref{refine:child-distance}, together
with $n_Q(d)\leq n_Q(b_j)$, give $d\leq\psi_Q(b_j)$.
Their minimum is therefore still an upper bound.  The argument holds
simultaneously for every solution satisfying the tangent constraint;
it does not select a single solution at any step.  Enlargement and
early stopping preserve the conclusion.
\end{proof}

The routine \texttt{relationB::eval} implements this argument with at
most 30 attempted improvements.  It accepts $0\leq t\leq2$ and otherwise
returns the non-success sentinel $10$, which exceeds every radius used
here.  On the accepted path it first enlarges $t$ by $2\cdot10^{-8}$,
adds $10^{-9}$ to the parent tangent factor, and uses $10^{-8}$-scale
enlargements in the initial distance, normal, child-tangent and updated
distance bounds.  A proposed value that does not improve the current
upper bound ends the iteration.  Section~\ref{sec:rounding} verifies
these enlargements, the bounds computed in \texttt{init}, and the final
$10^{-8}$ addition against the separate binary64 operations.  Denote
the resulting computed function by $\widehat D(t)$.  Its required
semantic property is the conclusion of
Lemma~\ref{refine:iteration-certificate}:
\begin{equation}\label{refine:eval-semantics}
 \|\eta-\eta_q\|\leq t
 \quad\Longrightarrow\quad
 \|\delta-d_q\|\leq\widehat D(t)
 \qquad(0\leq t\leq2).
\end{equation}
The proof of this property consists of the preceding induction plus
the explicit arithmetic bounds; the numerical snapping or
eigen-iterations are not premises of that induction.

There are two ways to use \eqref{refine:eval-semantics}.  One may
evaluate it at the farthest tangent distance of a particular box.
Alternatively, for a queried sphere radius $0\leq a\leq .126$, \texttt{init} searches
for a useful tangent-distance threshold $\rho$.  The search begins
with $\ell=0$ and $u=1.01a$, rejects if
$\widehat D(0)>a-10^{-7}$, and makes 25 midpoint trials.  A trial $m$
replaces $\ell$ precisely when
\begin{equation}\label{refine:rho-trial}
 \widehat D(m)<a-2\cdot10^{-7};
\end{equation}
otherwise it replaces $u$.  Finally the code sets
\begin{equation}\label{refine:rho-return}
 \rho=\max(0,\ell-10^{-8})
\end{equation}
and reports a usable relation only if $\rho>10^{-7}$.

\begin{lemma}\label{refine:rho-certificate}
Every usable returned threshold certifies
\[
 \|\eta-\eta_q\|\leq\rho
 \quad\Longrightarrow\quad
 [e(P+\delta)]\in\mathcal S(Q,a).
\]
The search need not find the largest possible threshold.  In particular,
its correctness does not require monotonicity of the floating-point
function $\widehat D$.
\end{lemma}

\begin{proof}
A positive returned $\rho>10^{-7}$ cannot arise from the initial value
$\ell=0$.  Thus the final stored $\ell$ was assigned by a successful
trial \eqref{refine:rho-trial}.  The arithmetic margins ensure
$\rho\leq\ell$ and $\widehat D(\ell)<a$.  Apply
\eqref{refine:eval-semantics} with this one stored trial value $\ell$.
It is a uniform certificate for every true tangent distance at most
$\ell$, and hence for every distance at most $\rho$.  The alignment
\eqref{refine:alignment} then gives the stated containment.  No
comparison between $\widehat D$ at distinct arguments is needed.
\end{proof}

For a previously completed parent used as a blocker, $a$ is its fixed
input radius and this lemma suffices.  For an atlas child, the search
uses the allowed cap $a=\tau$, not merely its current radius.  When
boxes are assigned to that child, the code evaluates $\widehat D$
again at an upper tangent distance for those boxes and grows the stored
radius to that certified distance, with an additional margin.  It
checks both that the tangent argument is at most $\rho$ and that the
new radius is at most $\tau$.  Failure of either check aborts.  Thus
even an unfavorable rounded reevaluation cannot silently assign a box
to a sphere too small to contain it.

\subsection{An exhaustive grid in tangent coordinates}

For every solution in the parent record,
\begin{equation}\label{refine:tangent-range}
 \|\eta\|=\|T^*\delta\|
       \leq\sqrt{1+o_P}\,R=:E_P.
\end{equation}
Let $w>0$ be the stored refinement width, regarded as an exact real
number.  The code forms an inflated upper bound $E$ for $E_P$ by adding
$10^{-7}$ to the computed right side of \eqref{refine:tangent-range},
and chooses
\[
 L=\left\lceil E/w+0.500001\right\rceil,
 \qquad n=2L+1,
\]
with the stated rounded evaluation.  The error bounds in
Section~\ref{sec:rounding} ensure that the resulting outer endpoints
$\pm(L+1/2)w$ contain $[-E_P,E_P]$.

For $u\in\R^k$ and $h\geq0$, write
\begin{equation}\label{refine:box-definition}
 B(u,h)=\prod_{a=1}^k[u_a-h,u_a+h].
\end{equation}
The base boxes have $u_a=b_aw$, $-L\leq b_a\leq L$, and $h=w/2$.
Their union contains $[-E_P,E_P]^k$, so every tangent vector in
\eqref{refine:tangent-range} belongs to at least one base box.  All
intervals are closed; grid boundaries are included.

The minimum squared distance of a box from the origin is
\begin{equation}\label{refine:outside-formula}
 \operatorname{dist}(0,B(u,h))^2
 =\sum_{a=1}^k\max(0,|u_a|-h)^2.
\end{equation}
The routine \texttt{outside} uses a lower version of this expression,
subtracting $10^{-8}$ inside each nonnegative part, and compares with
the inflated $E^2$.  Hence an accepted outside test means that the box
cannot contain the tangent vector of a parent solution.  Stored centers
are approximations to the exact multiples of $w$; repeated subdivision
centers are interpreted as exact dyadic offsets from this same grid.
The coordinate and arithmetic errors, bounded in
Section~\ref{sec:rounding}, are covered by these margins.

Subdivision is exhaustive because
\begin{equation}\label{refine:subdivision}
 B(u,h)=\bigcup_{\epsilon\in\{-1,1\}^k}
         B(u+\epsilon h/2,h/2).
\end{equation}
The code visits all $2^k$ children if they need to be examined; their
common boundaries remain covered.  When $k=0$, the empty product
\eqref{refine:box-definition} is the single point of $\R^0$, so there
is one base box and no subdivision is needed or permitted.

\begin{lemma}\label{refine:farthest-corner}
For every tangent box,
\begin{equation}\label{refine:corner-formula}
 \max_{\eta\in B(u,h)}\|\eta-\eta_q\|
 =\left(\sum_{a=1}^k(|u_a-\eta_{q,a}|+h)^2\right)^{1/2}.
\end{equation}
Consequently, if a certified upper bound $t_B\in[0,2]$ for this expression
satisfies $\widehat D(t_B)\leq r$, all parent solutions with tangent
vector in this box belong to $\mathcal S(Q,r)$.  The same statement
holds with $r=a$ if $t_B\leq\rho$ for a usable searched relation.
\end{lemma}

\begin{proof}
In each coordinate the maximum of $|\eta_a-\eta_{q,a}|$ on the
corresponding interval is $|u_a-\eta_{q,a}|+h$.  The endpoints attaining
these coordinate maxima can be chosen independently, so the maximum of
the sum of their squares is the sum in \eqref{refine:corner-formula}.
The conclusions follow from \eqref{refine:eval-semantics} and
Lemma~\ref{refine:rho-certificate}.
\end{proof}

The routine \texttt{corner} adds $10^{-8}$ after the computed norm to
obtain such an upper $t_B$.  The multi-box routine \texttt{markRec}
uses the same geometric fact on entire base boxes.  At each coordinate
it enlarges $|b_aw-\eta_{q,a}|+w/2$ by $10^{-8}$, accumulates its
square, and keeps a branch only while the partial sum is at most
$(\rho-5\cdot10^{-8})^2$.  After a complete box is marked, the bound
used for radius growth is the accumulated norm plus $2\cdot10^{-8}$,
or the maximum of these bounds when several boxes are marked.  The
coordinate inflations make this a true farthest-distance upper, while
the gap between $5\cdot10^{-8}$ and $2\cdot10^{-8}$ keeps it below
$\rho$ after rounding.  At $k=0$ the marked bound is $2\cdot10^{-8}$,
also smaller than every usable $\rho$.  The detailed error comparison
is given in Section~\ref{sec:rounding}.

In particular, a relation discovered while solving a proper subbox
does not justify marking its base box merely because that subbox is
covered.  A mark is made only after the independent whole-base-box
test just described.  Integer coordinate ranges and hash neighborhoods
only limit which marking or reuse attempts are made.  Any unmarked
box remains in the ordinary exhaustive traversal.

\subsection{Assignments, symmetry maps, and radius growth}

Let $\tau$ be the common cap on child radii.  The atlas stores a list
of certified Roots $Q_i$ with current radii $r_i\leq\tau$.  Its centers,
tangent columns and certificate constants remain fixed after insertion;
only $r_i$ may grow.

For a live box, the program may propose a new center by snapping a
matrix near $P+Tu$.  This is only a method of proposing a Root.  It
does not assert that $P+Tu$ is exact, that snapping converges, or that
the tangent coordinates describe a smooth locus.  A proposal must pass
the independent Root certificate and the parent--child relation test.
If $t_B$ is the box's corner bound, the direct path uses
$r=\widehat D(t_B)+2\cdot10^{-8}$ and requires
$r<\tau-10^{-7}$ before assignment.

\begin{lemma}\label{refine:atlas-assignment}
Suppose a box has been certified into $\mathcal S(Q,r)$, with $r<\tau$.
A signed permutation/transpose map $\Phi$ may be used to arrange $Q$.
If an existing atlas center $Z$ has a certified aligned distance
$\ell$ from $\Phi Q$, then growing its radius to any certified upper
bound for $r+\ell$ covers the box, provided that bound is at most
$\tau$.  Alternatively, inserting $\Phi Q$ with radius $r$ covers
the box.  A later failure to construct a reuse relation does not undo
either assignment.
\end{lemma}

\begin{proof}
For a solution in $\mathcal S(Q,r)$, take its horizontal displacement
$h$ with $\|h\|\leq r$.  Certificate transport and the commutation
$\Phi\Pi=\Pi\Phi$ give a horizontal displacement $\Phi h$ of the
same norm at $\Phi Q$.  By the definition of aligned distance, there
is a horizontal vector $a$ with $\|a\|\leq\ell$ such that $Z+a$ is
equivalent to $\Phi Q$ by integer phases and rephasings.  Therefore
$Z+a+\Phi h$ is an equivalent exact solution, and
$\|a+\Phi h\|\leq\ell+r$.  This proves the merge claim; the insertion
claim is the special case of retaining the transformed center itself.
Both are established before any later reuse attempt.
\end{proof}

The function \texttt{assign} implements this triangle argument.  For
a merge it requires $r+\ell+2\cdot10^{-8}<\tau$ and grows the old
radius to the maximum of its present value and
$r+\ell+10^{-8}$.  These margins are certified in
Section~\ref{sec:rounding}.  If no merge succeeds, it inserts the
canonical transformed Root with radius $r$.  The returned map is the
one used in the successful merge, or the canonical map in the insertion
case; no assertion that the search found the nearest atlas center is
needed.

For subsequent reuse, \texttt{inverseTransform} pulls the selected
atlas Root $Z$ back to $\Phi^{-1}Z$.  This is a signed
permutation of its stored phase and tangent coefficients, with its
constants unchanged, as justified by Lemma~\ref{phase:transport}.
The pulled-back Root and $Z$ have isometric solution spheres.  A fresh
relation to $\Phi^{-1}Z$ may therefore grow the radius stored at $Z$.
Any integer or rephasing alignment needed for that fresh relation is
chosen anew in \eqref{refine:alignment}.  A failure of this optional
\texttt{addNear} step cannot invalidate the already established direct
assignment.

Preloading also permits a map to be applied to the \emph{parent}.
This does not change the coordinates labeling its boxes.  Indeed, for
$V=\Phi P$, the transported tangent matrix and displacement satisfy
\begin{equation}\label{refine:mapped-coordinates}
 (\Phi T)^*(\Phi\delta)=T^*\delta=\eta.
\end{equation}
Thus a relation from the parent Root at $V$ to a stored child or blocker
can be used on the same grid indices.  The represented matrices change
only within their permitted equivalence classes.  Trying a limited
collection of maps, omitting hash neighbors, or applying an early
distance cutoff can only omit possible reuse certificates; none of
these operations declares a box covered by itself.

\subsection{Successful parents and the blocker induction}

A blocker is an earlier input parent whose entire refinement has
succeeded.  It is queried at its original fixed radius.  In contrast,
an atlas child is queried directly as a certified-center sphere and its
radius is grown when needed.  This distinction is important: a child
created during an unfinished parent search can already be reused,
whereas that unfinished parent cannot be used as a blocker.

\begin{proposition}\label{refine:cover-preservation}
Suppose Stage B is run on a finite list of records with valid Root
certificates, under the arithmetic assumptions and error bounds of
Section~\ref{sec:rounding}.  On a completed run, let $\mathcal A$ be its final atlas
of children and let $\mathcal B$ be the list of input parents whose
refinement failed.  Then
\begin{equation}\label{refine:cover-inclusion}
 \begin{split}
 \bigcup_{(P,R)\,\mathrm{input}}\mathcal S(P,R)
 &\subseteq
 \bigcup_{(Q,r)\in\mathcal A}\mathcal S(Q,r)\\
 &\qquad\cup
 \bigcup_{(P,R)\in\mathcal B}\mathcal S(P,R).
 \end{split}
\end{equation}
In particular, a completed run with no bad parents replaces the input
cover by a child cover.  A nonempty bad list remains an obligation.
\end{proposition}

\begin{proof}
Process the parents in the order used by the program.  The induction
invariant is that every previously successful parent sphere is covered
by the current atlas.  It holds before the first parent.  Appending
children and increasing radii preserve it.

For the current parent, first consider an individual box in the finite
subdivision tree.  A successful outside test proves that it contains
no parent solution.  A successful direct assignment covers all its
solutions by Lemmas~\ref{refine:farthest-corner} and
\ref{refine:atlas-assignment}.  A successful relation to a child covers
them by that child's checked radius update.  A successful relation to
a blocker covers them by the blocker sphere, which is already covered
by the atlas by the induction invariant.  Finally, success on all
children covers the box by \eqref{refine:subdivision}.  These are all
the successful returns of the cell solver, so induction from the
leaves proves its assertion for every successful cell.

The grid covers every possible tangent vector by
\eqref{refine:tangent-range}.  The traversal can finish successfully
only after every base box has been either rejected as outside, marked
by a valid whole-box relation, or solved as above.  Hence success for
the current parent proves coverage of its entire sphere and permits
its addition to the blocker list.  If the parent fails, it is not added
to that list and is retained in $\mathcal B$.  Any children already
produced while attempting it remain valid stored spheres; reusing them
does not assume that this failed parent has been covered.

The invariant therefore holds for every successful parent.  Failed
parents appear explicitly on the right of
\eqref{refine:cover-inclusion}, proving the claimed inclusion.
\end{proof}

This proof involves two finite inductions: one on subdivision depth and
one on the order of completed parents.  It does not infer coverage of a
parent from children that merely assume coverage of that same parent.
Moreover, later radius growth cannot invalidate an earlier certificate,
because centers and Root constants are immutable and the corresponding
solution sets only enlarge.

\subsection{Batch files and finite failure conditions}

Batching writes appended atlas Roots in their insertion order, in
groups of at most 256.  The radius written for each Root is the common
cap $\tau$, rather than its current achieved radius.  If a Root at
index $i$ is emitted before the run ends, every later state satisfies
\begin{equation}\label{refine:batch-containment}
 r_i\leq\tau,
 \qquad\mathcal S(Q_i,r_i)\subseteq\mathcal S(Q_i,\tau),
\end{equation}
with exactly the same center, tangent columns and certificate constants
as in its batch.  No future merge changes these data.  At successful
completion, the final partial batch and the count marker account for
every appended index.  Consequently replacing the child union in
\eqref{refine:cover-inclusion} by the union of all batch spheres at
radius $\tau$ is valid.

Each batch is completed under a temporary filename before atomic
rename.  Downstream computation may read these completed files while
Stage B continues, but the final conclusion still requires the complete
batch collection and the final Stage B failure checks.  The existence
of an intermediate batch, or even a count marker without the exit and
bad-list checks, is not a certificate that all parents were covered.
The exact execution protocol is specified in
Section~\ref{sec:execution}.

\begin{corollary}\label{refine:batch-cover}
Under the hypotheses of Proposition~\ref{refine:cover-preservation},
suppose a batching invocation completes with no bad parents and every
emitted batch is included.  Then the input solution sets are covered,
up to the permitted equivalences, by the emitted records with their
stored target radius $\tau$.
\end{corollary}

\begin{proof}
Use \eqref{refine:cover-inclusion} with the bad union empty, and enlarge
each final child sphere according to \eqref{refine:batch-containment}.
The immutable insertion order and the completed batch collection
include each final child center.
\end{proof}

For clarity, the implementation uses the following finite limits.
Input radii satisfy $0\leq R\leq0.126$, the child cap and width satisfy
\[
 0.0001\leq\tau\leq0.126,
 \qquad 0.0001\leq w\leq0.05,
 \qquad 0\leq k\leq10.
\]
The normal constants are nonnegative and are checked against
\[
 \gamma_P<20,\quad
 \mathtt{n0}_P,\varepsilon_P<10^{-6},\quad
 \alpha_P<10^{-4},\quad
 o_P,\mathtt{hor}_P<10^{-5},\quad
 \mu_P<2.
\]
Stored phase entries have magnitude at most $4$, and stored tangent
entries at most $1.01$.  These range checks accompany genuine Root
certificates; the range checks alone do not establish those
certificates.  A prospective Root with more than ten tangent columns
is rejected, not truncated.

The base grid is limited to $50{,}000{,}000$ boxes, subdivision depth is
at most six, and the atlas is limited to $2{,}000{,}000$ Roots.  A single
Record file is limited to $2{,}000{,}000$ records and the combined input
to $8{,}000{,}000$.  The grid product is checked before multiplication
past its limit.  Within the admitted ranges one has $L\leq1262$;
the mixed-radix indices fit the stated integer types, and every split
uses at most $2^{10}$ children.  The integer and floating-point bounds
are recorded in Section~\ref{sec:rounding}.

Exceeding the grid limit fails the parent.  An unresolved cell at the
depth limit, or an unresolved cell with $k=0$, also fails it.  Invalid
input, impossible radius updates, exhausted atlas capacity and file
errors abort the run.  An unsuccessful proposal or relation can lead
to subdivision or failure, but never directly to a discard.  In
particular, the clipped alignment tests in \texttt{relationB::init}
only decline a relation: their chosen cut is at most $0.249$ in the
direct-proposal path and at most $0.23$ in ordinary reuse.  They are
not exclusions of parent solutions.

The invocation parameters used for the three prescribed refinement
passes are respectively
\[
 (\tau,w)=(0.07,0.013),\qquad
 (0.016,0.005),\qquad
 (0.011,0.0038).
\]
The second pass uses batching; the third is applied to the records
retained by the first exclusion pass.  These choices affect whether
the finite search succeeds, but not
Proposition~\ref{refine:cover-preservation}.  The proposition is an
algorithmic implication and does not assume any claimed run counts,
hashes, or generated data.  Establishing that its successful-completion
hypothesis was met for the required records is a separate execution
obligation, addressed in Section~\ref{sec:execution}.
